\section{Marco Te\'{o}rico y Revisi\'{o}n de Literatura}
\label{sec:literature}

\subsection{Gesti\'{o}n de Inventarios}
\label{subsec:inventory}

[Fundamentos te\'{o}ricos: modelos cl\'{a}sicos (EOQ, modelo de Wilson, pol\'{i}ticas $(s, Q)$ y $(s, S)$), costos de inventario (ordenar, mantener, escasez), clasificaci\'{o}n ABC.]

\subsection{Anal\'{i}tica Prescriptiva}
\label{subsec:prescriptive}

[Definici\'{o}n y diferencias con anal\'{i}tica descriptiva y predictiva. Framework de \cite{bertsimas2020predictive}: de la predicci\'{o}n a la prescripci\'{o}n.]

\subsection{Machine Learning para Pron\'{o}stico de Demanda}
\label{subsec:ml-forecasting}

[Revisi\'{o}n de algoritmos aplicados: Random Forest, Gradient Boosting \parencite{chen2016xgboost, ke2017lightgbm}, SVR, redes neuronales. Comparaci\'{o}n con m\'{e}todos estad\'{i}sticos cl\'{a}sicos \parencite{carbonneau2008application}.]

\subsection{Sector Log\'{i}stico en el Per\'{u}}
\label{subsec:peru-logistics}

[Contexto del sector log\'{i}stico peruano: tama\~{n}o, desaf\'{i}os, nivel de adopci\'{o}n tecnol\'{o}gica, oportunidades de mejora.]

\subsection{Trabajos Relacionados}
\label{subsec:related-work}

[Estudios previos que aplican ML a inventarios. Gaps identificados en la literatura que esta tesis aborda.]
